\chapter{A motivating example}

% Goal: Reader should just be able to look at this chapter and know how Fulminate works at a high level.

% In this chapter, we consider a typical singly-linked list implemented in C and one of its operations, demonstrating how one can use Fulminate to gradually specify and verify the latter. 

In this chapter, we take a look at how Fulminate works at a high level, using a classic separation logic example of a singly-linked list, with the list and its operations implemented in C and its specifications written using CN. TODO: rewrite this para once chapter done

% {\color{violet}
% New section outline (?):

% \begin{itemize}
%   \item Start with the C you want to specify/test
%   \item Then go to separation-logic predicate you'd like to write
%   \item Then show CN
%   \item Then the rest, executing in Fulm, writing spec in CN and testing in Fulm. Is that where we show the Fulm output? Need to show the pre and postcondition code and where injection is happening, but also the actual \texttt{IntList} predicate in C and how it just easily translates to that
% \end{itemize}
% }

% \section{A singly linked-list and its operations, in C}

% \begin{center}
% \lstinputlisting[language=CNC]{code_snippets/append.c}
% \end{center}

% TODO: explain linked-list struct and imperative (non-recursive) implementation -- idiomatic C. Segue into needing \texttt{IntList} predicate to specify this.

\section{A singly-linked list, in C}\label{sec:linked-list-in-c}

Consider the following struct definition in C, representing a singly-linked list: 

\begin{center}
\lstinputlisting[language=C, numbers=none]{code_snippets/linked_list_def.h}
\end{center}

A \texttt{struct node} has, by this definition, a fixed-width integer value associated with it and a pointer to the next node in the list, also of type \texttt{struct node}. This definition can be unfolded recursively via the \texttt{next} member until this pointer is null, meaning the end of the linked list has been reached.

For example, the following C code uses this definition to create a linked list with elements 1, 2 and 3, starting at \texttt{i1}:

\begin{center}
\lstinputlisting[language=C, numbers=none, firstline=3]{code_snippets/linked_list_concrete.c}
\end{center}

This construction is done in reverse order, starting at the last element, such that the list can be ``chained'' together in memory as it is being constructed via the \texttt{next} pointers.

Now consider the following \texttt{list\_append} operation on two such singly-linked lists, also implemented in C:

\begin{center}
\lstinputlisting[language=C, firstline=3]{code_snippets/append.c}
\end{center}

This takes two arguments of type \texttt{struct node}, \texttt{xs} and \texttt{ys}. If \texttt{xs} is the null pointer, the function returns \texttt{ys}; otherwise, it traverses \texttt{xs} through to its last element, and sets the \texttt{next} pointer of that element to be \texttt{ys}.

To get higher assurance that this code behaves the way we expect it to, we might like to test it, or verify it, or both; this might naturally lead us towards wanting to write a specification for this function. 

\section{Specification desiderata}\label{sec:spec-desiderata}

Ideally, we would like our specification for \texttt{list\_append} to:

\begin{enumerate}
  \item Describe certain \emph{functional} properties -- e.g.~what the integer elements of the returned, concatenated list are in relation to the elements of the input lists.
  \item Describe \emph{ownership} properties, since we are dealing with low-level C code that involves explicit pointers and their fine-grained manipulation -- here, we care especially about the structure of the lists in memory on entry to and exit from the function.
  \item Be \emph{executable} for the purpose of testing, but also usable for verification if desired. The executability of the specification is important not just for standalone testing of the C function, but also for being able to test the specification itself as we write it in an incremental fashion.
\end{enumerate}

The first goal is achievable using Hoare logic, but the second is not -- for tractable ownership reasoning, separation logic is needed. Thus, for the third goal to also hold, the separation-logic specification we write for \texttt{list\_append} must be executable -- and this is not obviously possible on first glance, as we explain in the next section. 

\section{Executing separation-logic specifications}

\subsection{Standard separation logic is not executable}

To understand why this is the case, consider the following separation-logic predicate \texttt{IntList}, which formally describes a singly-linked list of the form we saw in Section~\ref{sec:linked-list-in-c}:

\begin{figure}[h]
  \vspace{-0.8em}
  \begin{align*}
  \text{IntList}(p, \alpha) =\, &(p = \text{NULL} \land \alpha = \lbrack\rbrack)\\
  &\lor (\exists{hd}, q, {tl}.\,p \mapsto ({hd}, q) \ast \text{IntList}(q, {tl}) \land \alpha = {hd} :: {tl})
  \end{align*}
  \vspace{-2em}
\end{figure}

This predicate takes two inputs: a pointer $p$ that is either the null pointer, or points to the head of some linked list in memory, and a mathematical list $\alpha$, an abstract representation of the linked list's values. In the case where $p$ is not null, this predicate asserts ownership over the entire list one node at a time, via the points-to assertion and the recursive call to \texttt{IntList}.

What would it mean to execute this predicate as part of some function-level specification? Rather than proving that these properties hold statically for any arbitrary $p$ and $\alpha$, we want to check for some concrete values of $p$ and $\alpha$, and the corresponding concrete heap, that we can assert the correct ownership over all nodes of the linked list, and that all of the expected logical properties hold. 

However, this predicate does not look particularly executable on first glance. There is a disjunction, which might require backtracking depending on whether $p$ is the null-pointer or not. In this case, this is detectable early in the execution of the first disjunct, but for arbitrary disjunction, this is not guaranteed. In addition, the second disjunct contains existential quantification over ${hd}$, $q$ and ${tl}$, which might require searching for values, Furthermore, there is an instance of the separating conjunction $\ast$, which requires searching for disjoint splittings of the heap -- this blows up quickly with the size of the heap, with a heap of size $n$ requiring $2^n$ possible splittings to be searched.

Thus, executing a classical separation-logic specification is, in the general case, computationally infeasible. The next section describes how carefully restricting the logic can eliminate these barriers to executability.


\begin{figure}[h]
  \begin{center}
    \begin{minipage}[t]{.31\textwidth}
      \vspace{-6.2em}
      \lstset{
      numbers=left,
      firstnumber=1,
      numberfirstline=true
      }
      % \vspace{0.2em}
      \scalebox{0.9}{\lstinputlisting[language=CNC, lastline=12]{code_snippets/intlist.cn.c}}
    \end{minipage}
    %
    %
    \hspace*{8.5\baselineskip}
    \begin{minipage}[t]{.28\textwidth}
      \lstset{
      numbers=left,
      firstnumber=13,
      numberfirstline=true
      }
      \scalebox{0.9}{\lstinputlisting[language=CNC,firstline=13]{code_snippets/intlist.cn.c}}
    \end{minipage}
    % \vspace*{0.8\baselineskip}
  \end{center}
  \vspace*{-4mm}
  \caption{The \texttt{IntList} predicate, written in the CN specification language. The CN definitions are enclosed in magic comments {\color{commentcolour}$/*@ \ldots @*/$} (lines 1-12, 18-23). The \texttt{IntList} predicate definition relies on the the C \texttt{struct node} definition (lines 13-16), which is the same as the one introduced in Section~\ref{sec:linked-list-in-c}.}
  \label{fig:intlist-cn}
\end{figure}

\subsection{CN separation logic \underline{is} executable!}

Now consider the same \texttt{IntList} predicate, written in the CN specification language, as shown in Fig.~\ref{fig:intlist-cn}. The CN predicate definition is on the left hand side, with auxilliary C and CN definitions on the right. The CN definitions are enclosed in ``magic comments'' ({\color{commentcolour}$/*@ \ldots @*/$}), such that they can be parsed and desugared by the CN toolchain and treated like usual comments by any C compiler such as clang or gcc. 

\TODO{rename node variable to something else? it's confusing for the variable and the struct tag to have the same name}

The \texttt{IntList} predicate signature on line 2 states that the predicate takes a single argument $p$ of generic pointer type, and returns a \texttt{datatype seq}, which is defined on lines 19-22 as a CN datatype with two constructors, \texttt{Nil} and \texttt{Cons}. The \texttt{rec} keyword indicates that the predicate is recursive. The predicate first checks whether the pointer $p$ is null, returning the \texttt{Nil} constructor if that is the case (lines 3-5). Otherwise, it extracts the value stored at $p$ using CN's builtin \texttt{Owned} predicate, storing it in a new variable called \texttt{node} (line 7). (Note the angle brackets $\langle\ldots\rangle$, which denote the expected C type of $p$'s value -- here, this has type \texttt{struct node}, defined on lines 13-16 as a usual C singly-linked list with some value \texttt{val} and a \texttt{next} pointer to the rest of the list.) After that, \texttt{IntList} is called recursively on the \texttt{next} pointer of $p$'s value, with the associated mathematical list datatype stored in a new variable \texttt{tl} (line 8). Finally, the 32-bit integer \texttt{head} member of \texttt{node} (which we recall has type \texttt{struct node}) and the mathematical abstraction of the rest of the list are composed under the \texttt{Cons} constructor and returned (line 9).  

\TODO{Following para might need to be rejigged slightly to improve flow}

This CN predicate differs from its classical analogue in several ways: it takes the pointer $p$ as its only argument rather than both the pointer and its mathematical list representation $\alpha$; it has an if-else statement that checks whether $p$ is NULL instead of using logical disjunction to account for either case; and it constructs and returns the mathematical list $\alpha$ as a CN datatype. Crucially, points-to assertions are written in CN using its builtin \texttt{Owned} predicate: $p \mapsto v$ is written as $\texttt{take}\,\,v = \texttt{Owned}\langle\textit{ctype}\rangle(p)$, with CN's syntax restricting $v$ to be an output parameter here rather than an input that could be arbitrarily quantified over. The $\ast$ operator is encoded by CN's semicolons, ensuring each assertion holds for disjoint parts of the heap.

This version of the \texttt{IntList} predicate looks executable, and it turns out that it is. CN's restrictive syntax, originally designed to make the CN proof tool's automation more predictable, is exactly what enables it to be executed -- and this is what Fulminate sets out to do.

\subsection{Fulminate executes CN specifications}

Fulminate is a tool that takes a C file with CN annotations, and transforms these into C runtime checks, producing an instrumented C file with executable separation-logic specification checks. For example, Fulminate takes the CN \texttt{IntList} example in Fig.~\ref{fig:intlist-cn} as input and produces the C code in Figs~\ref{fig:intlist-fulm-datatype-defs} and \ref{fig:intlist-fulm}.

\begin{figure}[h]
  \begin{center}
    \begin{minipage}[t]{.15\textwidth}
      % \vspace{-6.2em}
      \lstset{
      numbers=left,
      firstnumber=1,
      numberfirstline=true
      }
      % \vspace{0.2em}
      \scalebox{0.9}{\lstinputlisting[language=CNC, firstline=5, lastline=20]{code_snippets/intlist.cn.exec.c}}
    \end{minipage}
    %
    %
    \hspace*{8.5\baselineskip}
    \begin{minipage}[t]{.3\textwidth}
      \lstset{
      numbers=left,
      firstnumber=17,
      numberfirstline=true
      }
      \scalebox{0.9}{\lstinputlisting[language=CNC,firstline=22, lastline=38]{code_snippets/intlist.cn.exec.c}}
    \end{minipage}
    % \vspace*{0.8\baselineskip}
  \end{center}
  \vspace*{-4mm}
  \caption{Top-level type definitions produced by Fulminate when applied to CN \texttt{IntList} example.}
  \label{fig:intlist-fulm-datatype-defs}
\end{figure}

\begin{figure}[h]
  \begin{center}
  \lstset{
    numbers=left,
    firstnumber=34,
    numberfirstline=true
  }
  \begin{minipage}[t]{\textwidth}
  \scalebox{0.9}{\lstinputlisting[language=C, firstline=41, lastline=58, breaklines]{code_snippets/intlist.cn.exec.c}}
  \end{minipage}
  \end{center}
  \caption{Fulminate translation of the CN \texttt{IntList} predicate.}
  \label{fig:intlist-fulm}
\end{figure}

Fig.~\ref{fig:intlist-fulm-datatype-defs} shows the top-level type definitions in the output file. Fulminate preserves the \texttt{node} struct definition from the source, and produces a mirrored definition of it with tag \texttt{node\_cn} (lines 7-10), whose members' types belong to the set of Fulminate's runtime representation of CN types. Here, the C signed int from line 2 is mapped to CN's signed 32-bit bitvector type on line 8, which Fulminate represents under the hood as a boxed C \texttt{int32\_t}. Similarly, the \texttt{next} member's type is mapped from a pointer to a \texttt{struct node} into a generic CN pointer, which Fulminate represents as a boxed \texttt{void *} pointer. 

Fulminate translates the CN \texttt{seq} datatype to the \texttt{seq} struct definition on lines 29-32, which in turn relies on the type definitions on lines 13-27. Every \texttt{struct seq} in the output carries a tag for discerning which constructor is being used in a given instantiation of the struct, and a struct representing the constructor's members and values, which is wrapped in a bespoke union. For example, a \texttt{Nil \{\}} in CN is represented as a \texttt{struct seq} in Fulminate, whose tag is \texttt{NIL} (line 14) and whose union member contains a pointer to some \texttt{struct nil} (line 17), which has no members, as expected.

Now we can inspect Fulminate's translation of the \texttt{IntList} predicate, as shown in Fig.~\ref{fig:intlist-fulm}. This is a C function that takes a pointer to some generic CN pointer $p$ as its first argument, analagously to the CN predicate definition. The Fulminate \texttt{IntList} function and some of its callees (namely, \texttt{owned\_struct\_node} and its recursive call to itself) take some additional auxilliary arguments as signified by the $\ldots$; we omit these details here.{\color{red}TODO: bring back spec\_mode argument so we can discuss it in 3.4? Keep \ldots for ignoring \texttt{loop\_ownership} argument} This function has the same shape as the CN predicate: the if-statement on line 37 checks if the pointer $p$ is null, returning Fulminate's representation of \texttt{Nil \{\}} if that is the case. Otherwise, Fulminate calls its bespoke \texttt{owned\_struct\_node} function on $p$, which first dynamically asserts ownership of $p$, using machinery we describe in Chapter~\ref{chap:impl}. If the assertion of ownership is successful, the C pointer associated with $p$ under the hood is dereferenced to get back a value of type \texttt{struct node}, which is then converted to its corresponding CN struct type and stored in the newly-created \texttt{node} variable. (Fulminate generates analagous functions named \texttt{owned\_\textit{ctype}} for all user-defined types \texttt{\textit{ctype}} in the source.) Then, the \texttt{IntList} function is called recursively on line 44, with the \texttt{next} member of the CN-mirrored struct \texttt{node} being passed to it (which we know to conveniently have type \texttt{cn\_pointer *}, as defined on line 9 and required by the signature of the generated \texttt{IntList}). Finally, Fulminate's representation of \texttt{Cons} is allocated for in memory (lines 45, 47) and assigned the expected values, using the \texttt{val} member of \texttt{node} (line 48) and the result of the recursive call to \texttt{IntList} (line 49), and then finally returned (line 50), just as in the CN \texttt{IntList} definition.

\bigskip

We can now see that introducing some careful restrictions to separation logic makes it possible to execute and test separation-logic specifications. CN's syntax is intentionally restrictive, disallowing arbitrary quantification and logical disjunction, and restricting arbitrary points-to assertions by assigning input and output modes to their pointer and pointee arguments respectively. Fulminate makes use of these restrictions, transforming CN separation-logic specifications into C executable specifications that check functional and ownership properties at runtime.

In this section, we have sketched out the separation-logic predicate \texttt{IntList} that we need to specify the singly-linked list introduced in in Section~\ref{sec:linked-list-in-c}, first in standard separation logic and then in CN's restricted form of it, and decided the CN toolchain is most amenable to our specification-writing goals (Section~\ref{sec:spec-desiderata}), due to its support for executable separation-logic specifications via Fulminate. Now we can return to the task at hand: writing a specification for the \texttt{list\_append} function.


% \section{Executing \texttt{IntList}}

\section{Writing a specification for \texttt{list\_append}, using Fulminate}

Recall the \texttt{list\_append} function we first saw in Section~\ref{sec:linked-list-in-c}:

\begin{figure}[h]
    \begin{center}
    \lstinputlisting[language=C, firstline=3]{code_snippets/append.c}
    \end{center}
\end{figure}


We can now use CN's specification language to start describing the arbitrary intended behaviour of the function. In particular, we expect that on entry to the function, \texttt{xs} and \texttt{ys} are either null, or pointers to distinct singly-linked lists in memory\footnote{something about the assumption of disjointness being a simplification for specification purposes?}, since both have type \texttt{struct node *}.

This leads us to our first specification attempt for \texttt{list\_append}, as shown in Fig.~\ref{fig:partial-cn-spec}. The \texttt{requires} keyword in CN signifies a function precondition. The first clause, \texttt{take l1 = IntList(xs)}, applies our \texttt{IntList} predicate to CN's representation of \texttt{xs}, which has CN's generic \texttt{pointer} type, and binds the output of the predicate to a new variable \texttt{l1}, which we know from the definition of \texttt{IntList} (Fig.~\ref{fig:intlist-cn}) to be of type \texttt{datatype seq}. This holds analagously for the second clause over \texttt{ys}.

\begin{figure}[h]
    \begin{center}
    \lstinputlisting[language=C, firstline=21]{code_snippets/append.cn.1.c}
    \end{center}
    \caption{An initial, partial CN specification for \texttt{list\_append}.}
    \label{fig:partial-cn-spec}
\end{figure}

Even at this early stage of specification-writing, before writing e.g.~a postcondition or any inline assertions, it would be useful to know whether our specification is on the right track by testing what we have so far, for which we can apply Fulminate to the partially-specified \texttt{list\_append }function.

\subsection{Fulminate's instrumentation of \texttt{list\_append}}

Applying Fulminate to the code in Fig.~\ref{fig:partial-cn-spec} yields the instrumented \texttt{list\_append} in Fig.~\ref{fig:partial-fulm-spec}, simplified slightly for presentation purposes. The purpose of each part of the instrumentation is described in detail below.

\emph{Initial setup on function entry.} TODO properly: yada yada, map every C function argument to its corresponding CN type and bind to a new variable to be used in all spec from that point onwards.

\emph{Function-level specifications.} For function-level specifications, Fulminate injects its executable C translation of each kind of specification into the expected place in the function. Namely, executable preconditions are injected to the start of the function, before the body, such that they are checked immediately on function entry. Postconditions require a bit more care: a \texttt{\_\_cn\_epilogue} label is injected at the end of the function, after its body, and the executable postcondition checks are placed after this label. Every return statement in the function is replaced with a block containing a statement that stores the value being returned at that point in a bespoke Fulminate-generated variable called \texttt{\_\_cn\_ret}, followed by a jump to the \texttt{\_\_cn\_epilogue} label to trigger the postcondition checks, after which the \texttt{\_\_cn\_ret} is returned. Between these two statements, some necessary ownership tracking statements may also be injected, which we will return to later.

In Fig.~\ref{fig:partial-fulm-spec}, the executable precondition corresponding to the CN precondition in the source is on lines 10 and 11. The executable postcondition instrumentation is generated regardless of whether a CN postcondition is provided or not; so far, we have not written a postcondition, so the statements under the \texttt{\_\_cn\_epilogue} label on line 35 do not correspond to any user-provided postcondition, but are, again, necessary for Fulminate's dynamic ownership tracking scheme to work. The two return statements from the source (lines 6 and 13 of Fig.~\ref{fig:partial-cn-spec}) are indeed transformed to blocks of statements that end in a jump to the executable postcondition (lines 18-19, 28-30 of Fig.~\ref{fig:partial-fulm-spec}), as described above.

\emph{Dynamic ownership tracking.}

\begin{figure}[h]
    \begin{center}
    \lstinputlisting[language=C, firstline=5]{code_snippets/append.cn.1.exec.c}
    \end{center}
    \caption{The partially-specified \texttt{list\_append} function from Fig.~\ref{fig:partial-cn-spec}, instrumented with the corresponding Fulminate runtime specification checks.}
    \label{fig:partial-fulm-spec}
\end{figure}


\newpage

% Outline

% \begin{itemize}
%   \item use \texttt{IntList} predicate to specify ownership, structure in memory (both encompassed by ``ownership pattern''?) and mathematical representation of list elements for xs and ys
%   \item use \texttt{requires} keyword, explain it's the keyword for a precondition in CN
%   \item show concrete test case -- \texttt{main} function needed, tiny handwritten driver. easy to write! emphasise ease
%   \item explain what application of Fulminate does abstractly. Might not need another figure, or maybe just comments/single function call for pre and post (like Saljuk actually has in Silver Fulm). Or put figure in appendix and just discuss abstractly here
%   \item show the runtime error and lldb trace after just providing precondition. (see overview/walkthrough section of PLDI paper). Emphasise that this is a partial spec. Fragmentary not so important here, unless you say the call to \texttt{list\_append} is wrapped inside some function \texttt{f} (might be too complex at this point, not sure)
%   \item use these to debug and realise you need the postcondition too. Fill it in, get rid of runtime error.
%   \item then finally get to functional property, that the returned mathematical list coming out of \texttt{IntList} is the mathematical representations of \texttt{xs} and \texttt{ys} appended together (bring in CN logical function \texttt{append} at this point and talk about its Fulminate representation being analagous to \texttt{IntList} but without the ownership assertion basically. No resources allowed! double-check CN paper for phrasing). Probably don't need to get this one wrong and then right; they already get the picture of how to debug partial specs.
% \end{itemize}



% We can now consider the following C implementation of a singly-linked list and its append operation, specified in CN using the \texttt{IntList} predicate we just defined:

% \TODO{fix C spacing}

% \TODO{Check this new imperative implementation actually works with CN and Fulminate, with same spec}

% \begin{figure}[h]
%   \begin{center}
%   \lstinputlisting[language=CNC]{code_snippets/append.cn.c}
%   \end{center}
% \end{figure}





% Do example with partial specifications. Include error messages, debugging with \texttt{lldb}, etc

% \section{The CN separation-logic specification language}

% Cover $p \mapsto v$ and \texttt{Owned}
% IntList example (should be in above spec example) and side-by-side comparison between standard SL formulae and CN representation. 


% \section{Writing partial CN specifications using Fulminate}

% Assuming we have a pre-existing C function to specify in this case.

% Small spec example, using IntList. Explain CN semicolon is $\ast$. 

% Say Fulminate takes CN+C files, gives C file with checks in the expected places. Fulminate tracks ownership of stack-local and heap addresses; instruments access checks to ensure ownership of an address has been asserted properly before it is read from or written to.

% Show Fulminate runtime ownership error from partial spec, debug and fix up spec, run again.


